\begin{textsample}{POS Dim 9 – global\_north\_2024\_01 – Score 1.00 – global\_north\_2024\_01/200148198.txt}  \label{ex:f9_pos_269}
African leaders criticize Israel ' s military campaign in Gaza and call for an immediate cease-fire African leaders have criticized Israel for its military campaign in Gaza and called for an end to the fighting that continues take its toll on mostly civilians ByRISDEL KASASIRA Associated Press January 19, 2024, 12:32 PM Dennis Francis, president of the U.N.
General Assembly, speaks during the 19th summit of Non-Aligned Movement ( NAM ) in Kampala, Uganda, Jan. 19, 2024.
President of the UN General Assembly Dennis Francis and African Union Commission Chairman Moussa Faki Mahamat called for an immediate ceasefire in Gaza. ( AP Photo/Hajarah Nalwadda ) The Associated Press KAMPALA, Uganda -- African leaders on Friday criticized Israel for its military campaign in Gaza and called for an end to the fighting that continues take its toll on mostly civilians.
The African Union Commission Chairman Moussa Faki Mahamat described the war in Gaza as immoral and unacceptable."We demand an immediate end to this unjust war against Palestinians and implementation of the two-state solution,"he said. @ @ @ @ @ @ @ @ @ @ have died, and the United Nations says a quarter of the 2.3 million people trapped in Gaza are starving.
In Israel, around 1,200 people were killed during the Oct. 7 attack by Hamas that sparked the war and saw some 250 people taken hostage by militants.
Mahamat was speaking at a conference in Kampala of the Non-Aligned Movement ( NAM ), a group of 120 states which aspire not to be formally aligned with or against any major power bloc.
Speaking during the meeting of heads of state at the week-long gathering, Mahamat asked the 120 member countries to demand international justice for the Palestinians.
His remarks were echoed by South African President Cyril Ramaphosa, who called for the release of all the hostages and"the resumption of talks on a just solution that will end the suffering of the Palestinian people."Ramaphosa further called for \textbf{unhindered} and expanded humanitarian access to allow for vital aid and basic services to meet the needs of everyone living in Gaza.
South Africa has filed @ @ @ @ @ @ @ @ @ @ for genocide and has asked the United Nations ' top court to order an immediate halt to Israeli military operations in Gaza."This is necessary to protect against further, severe and irreparable harm to the rights of the Palestinian people,"Ramaphosa said.
At the start of the conference on Monday, the Palestinian ambassador to the U.N. called on the members of the Non-Aligned Movement to put pressure on Israel to implement a cease-fire in Gaza after 100 days of war with Hamas.
In his opening speech, Ambassador Rayid Mansour said despite resolutions by the U.N.
General Assembly and the Security Council, a cease-fire remained elusive.
The Non-Aligned Movement, formed during the collapse of the colonial systems and at the height of the Cold War, has played a key part in decolonization processes, according to its website.
Mansour compared Israel ' s military assault on Gaza to apartheid, the system of white minority rule in South Africa which was finally abolished in 1994.
Israel rejects such allegations.

% matched lemmas: unhindered
\end{textsample}
